\section{Математические модели в экономике}
\label{sec:models-in-economics}

Экономическая модель выделяет существенные связи между ресурсами,
выпуском, ценами, спросом, предложением и целевыми показателями. В
зависимости от задачи используются балансовые, оптимизационные,
динамические, вероятностные и статистические модели.

\subsection{Межотраслевой баланс Леонтьева}

Пусть $x_i$ --- валовой выпуск отрасли $i$, $y_i$ --- конечный спрос, а
$a_{ij}\ge0$ --- количество продукции отрасли $i$, необходимое для
производства единицы продукции отрасли $j$. Тогда
\[
x_i=\sum_{j=1}^n a_{ij}x_j+y_i.
\]
В матричном виде
\[
\boxed{x=Ax+y},
\qquad
\boxed{(I-A)x=y}.
\]
Если спектральный радиус неотрицательной матрицы прямых затрат удовлетворяет
\[
\boxed{\rho(A)<1},
\]
то матрица $I-A$ обратима и
\[
(I-A)^{-1}=I+A+A^2+\cdots\ge0.
\]
Следовательно, для любого $y\ge0$ существует единственный неотрицательный
вектор выпуска
\[
\boxed{x=(I-A)^{-1}y}.
\]
Это один из основных результатов модели продуктивной экономики Леонтьева.

\subsection{Оптимизация производственного плана}

Пусть предприятие выпускает $n$ видов продукции в количествах
$x_1,\ldots,x_n$. Если $c_j$ --- прибыль с единицы продукции, а
$a_{ij}$ --- расход ресурса $i$ на единицу продукции $j$, то задача
линейного программирования имеет вид
\[
\boxed{
\max_{x\ge0}\;c^Tx
\quad\text{при}\quad
Ax\le b.
}
\]
Здесь $b_i$ --- доступный запас ресурса $i$. Ограничения задают множество
допустимых производственных планов, а целевая функция --- выбранный
экономический критерий.

В линейном программировании при существовании конечного оптимума он
достигается, в частности, в одной из крайних точек допустимого многогранника.
На этом основаны симплексные методы. Двойственные переменные можно
интерпретировать как предельные, или теневые, цены ресурсов.

\subsection{Транспортная задача}

Пусть $a_i$ --- запасы у поставщиков, $b_j$ --- потребности потребителей,
$c_{ij}$ --- стоимость перевозки единицы груза. Необходимо выбрать объёмы
$x_{ij}\ge0$:
\[
\boxed{
\min\sum_{i=1}^m\sum_{j=1}^n c_{ij}x_{ij}
}
\]
при условиях
\[
\sum_jx_{ij}=a_i,
\qquad
\sum_ix_{ij}=b_j.
\]
Сбалансированная задача удовлетворяет
\[
\sum_i a_i=\sum_j b_j.
\]
Если суммы не совпадают, вводят фиктивного поставщика или потребителя.
Транспортная задача является специальным случаем линейного
программирования и имеет более эффективные специализированные алгоритмы.

\subsection{Модель спроса и предложения}

Простейшие линейные зависимости можно записать как
\[
Q_d(p)=a-bp,
\qquad
Q_s(p)=c+dp,
\qquad b,d>0.
\]
Рыночное равновесие определяется условием
\[
Q_d(p^*)=Q_s(p^*),
\]
откуда
\[
\boxed{p^*=\frac{a-c}{b+d}},
\qquad
Q^*=Q_d(p^*)=Q_s(p^*).
\]
Изменение параметров спроса или предложения позволяет исследовать
сравнительную статику равновесия.

\subsection{Максимизация прибыли}

Если предприятие выбирает объём выпуска $q$, цена продукции равна $p$, а
издержки задаются функцией $C(q)$, то
\[
\Pi(q)=pq-C(q).
\]
Для внутреннего оптимума необходимое условие
\[
\boxed{p=C'(q^*)},
\]
то есть цена равна предельным издержкам. Если $C$ строго выпукла,
$C''(q)>0$, то такой стационарный объём является единственным максимумом
прибыли на внутренности допустимого множества.

\subsection{Динамические модели}

Экономические характеристики могут изменяться во времени. Например,
простейшая модель накопления капитала имеет вид
\[
\dot K=sF(K)-\delta K,
\]
где $F(K)$ --- производственная функция, $s$ --- норма сбережений,
$\delta$ --- коэффициент выбытия. Стационарные состояния удовлетворяют
\[
sF(K^*)=\delta K^*.
\]
Такие модели исследуют методами теории динамических систем: находят
равновесия, изучают их устойчивость и влияние параметров.

\subsection{Что важно сказать на экзамене}

Лучше показать несколько принципиально разных классов: \textbf{балансовая
модель Леонтьева}, \textbf{линейное программирование и транспортная задача},
\textbf{равновесие спроса и предложения} и простую динамическую модель.
Для Леонтьева полезно записать условие $\rho(A)<1$ и формулу
$(I-A)^{-1}=I+A+A^2+\cdots$. В оптимизационных моделях обязательно назвать
переменные, ограничения и целевую функцию.

\newpage
